%% 第 18 章：伴随的矩阵
%% 本文件由 tools/md2tex.py 自动生成，请勿直接编辑。

\chapter{伴随的矩阵}


\begin{jprop}{规范正交基下的共轭转置}{r:19:1}
若 $V,W$ 分别选取规范正交基，且：

\[
[T]=A,
\]

则：

\[
[T^*]=A^*.
\]

这里：

\[
A^*={A}^{\mathsf H}
\]

称为 $A$ 的共轭转置。

在实数域上：

\[
A^*=A^{\mathsf T}.
\]
\end{jprop}

\section{为什么必须使用规范正交基}

只有在规范正交基下，坐标中的内积才直接等于普通欧氏点积。

若基不是规范正交基，坐标向量之间的普通点积并不等于原空间中的内积。此时必须先通过 Gram 矩阵记录该基的内积关系，再进行相应转换。

因此，伴随矩阵等于共轭转置”是规范正交坐标下的结论
